\renewcommand{\Anno}{2005}

\renewcommand{\NumeroProblema}{3}

\renewcommand{\Testo}{La Scuola Sant'Anna ha deciso di partecipare al
campionato di Formula Uno con una vettura, e si sta preparando per il prossimo gran premio che si svolgerà sul circuito di Montecarlo (comune della provincia di Lucca).  I suoi ingegneri hanno a disposizione i seguenti dati:
	\begin{itemize}
		\item  la gara consiste in 120 giri di pista;
		
		\item la capacità del serbatoio della vettura permette, se necessario, di completare la gara senza alcun pit-stop;
	
		\item per un pit-stop si impiegano 30 secondi in totale, compresi i tempi di rallentamento ed accelerazione, indipendentemente dalla quantità di carburante immessa;
	
		\item per ogni giro si consumano, indipendentemente da ogni altro parametro, 3 kg di carburante;
		
		\item il tempo di percorrenza di un giro aumenta di 3 centesimi di secondo per ogni kg di carburante in più presente nel serbatotio all'inizio del giro stesso.
	\end{itemize}
	
Determinare la strategia (cioè il numero dei pit-stop da effettuare ed i giri nei quali effettuarli) che permette di terminare
la gara nel minor tempo possibile.}

\renewcommand{\Risposta}{Conviene fare 4 pit-stop equispaziati,
cioè al termine dei giri 24, 48, 72, 96.}

\renewcommand{\Aiutino}{Indichiamo preliminarmente con $T$ il tempo (in centesimi di secondo)
che serve per percorrere un giro iniziandolo con soli 3 kg di carburante.  Nella risoluzione si può procedere secondo i seguenti passi.
\begin{itemize}
	\item  Dimostrare che il tempo necessario per percorrere $k$ giri è maggiore od uguale di
	$$kT+\frac{9}{2}k(k-1),$$
	con uguaglianza se e solo se si arriva alla fine con il serbatoio scarico.

	\item Supponendo di suddividere la gara in $n$ tratti di lunghezza $k_{1},\ldots,k_{n}$, dimostrare che il tempo totale di percorrenza è maggiore od uguale di
	$$120\left(T-\frac{9}{2}\right)-3000+\frac{9}{2}
	\sum_{i=1}^{n}k_{i}^{2}+3000n,$$
	con uguaglianza se e solo se si arriva alla fine di ogni tratto con il serbatoio scarico. Non resta dunque che minimizzare la somma degli ultimi 2 termini con il vincolo che $k_{1}+\ldots+k_{n}=120$.

	\item Dimostrare che tale somma è maggiore od uguale di
	$$600\left(\frac{108}{n}+5n\right),$$
	con uguaglianza se e solo se i $k_{i}$ sono tutti uguali, a costo di non essere interi.

	\item  Dimostrare (mediante uno studio di funzione od una disuguaglianza elementare) che tale espressione, al variare di $n$ tra gli interi positivi, assume minimo per $n=5$.

	\item  Osservare che dividendo la gara in 5 tratti di ugual lunghezza (quindi 4 pit-stop) si hanno tutte uguaglianze, dunque la scelta ottimale.
\end{itemize}}

\renewcommand{\SoluzioneUfficiale}{Conveniamo di indicare tutti i 
tempi in centesimi di secondo, e sia $T$ il tempo che serve per percorrere un giro iniziandolo con 3 soli kg di carburante nel serbatotio (cioè il minimo che serve per completarlo). Dividiamo
ora la soluzione in vari passi.\medskip

\emph{Passo 1}. Determiniamo il tempo minimo necessario per percorrere $k$ giri consecutivi senza soste. È chiaro che tale tempo minimo si avrà se si imbarca il minimo di carburante necessario. In base alle informazioni date, sappiamo che il tempo di percorrenza diminuisce di 9 centesimi ogni giro (3 centesimi per ciascuno dei 3 kg di carburante che si consumano). Pertanto, se l'ultimo giro viene percorso in $T$ centesimi, il penultimo richiede $T+9$ centesimi, il terz'utlimo $T+18$ centesimi, e cos\`{\i} via. Ne segue che il tempo minimo per percorrere $k$ giri è
\begin{eqnarray}
	T+(T+9)+(T+18)+\ldots+(T+(k-1)) & = & kT+9(1+2+\ldots+(k-1))
	\nonumber  \\
	 & = & kT+9\frac{k(k-1)}{2}
	 \label{stimatempi1}
\end{eqnarray}
centesimi di secondo, con uguaglianza se e solo se si imbarca la quantità esatta di carburante.
\medskip

\emph{Passo 2}. Supponiamo di suddividere la gara in $n$ tratti di lunghezza $k_{1},\ldots,k_{n}$, con $k_{1}+\ldots+k_{n}=120$, compiendo quindi $n-1$ rifornimenti.

Applicando la (\ref{stimatempi1}) su ogni tratto, e tenendo conto che 
ogni rifornimento dura 3000 centesimi, il tempo totale sarà allora
\begin{eqnarray}
	\mbox{Tempo Totale} & \geq & \sum_{i=1}^{n}
	\left(k_{i}T+\frac{9}{2}k_{i}^{2}-\frac{9}{2}k_{i}\right)
	+3000(n-1) 
	\nonumber \\
	 & = & \left(T-\frac{9}{2}\right)\sum_{i=1}^{n}k_{i}+
	 \frac{9}{2}\sum_{i=1}^{n}k_{i}^{2}+3000(n-1) 
	 \nonumber \\
	 & = & \left(T-\frac{9}{2}\right)\cdot 120-3000+
	\frac{9}{2}\sum_{i=1}^{n}k_{i}^{2}+3000n,
	\label{stimatempi2}
\end{eqnarray}
con uguaglianza se e solo se la vettura arriva scarica al termine di ogni tratto.

I primi due addendi non dipendono dalla strategia: non ci resta 
dunque che minimizzare la somma degli ultimi due addendi.\medskip

\emph{Passo 3}. Per la disuguaglianza tra media quadratica e media aritmetica (valida per ogni $n$-upla di numeri reali $k_{1},\ldots,k_{n}$) si ha che
$$\left(\frac{1}{n}\sum_{i=1}^{n}k_{i}^{2}\right)^{1/2}\geq
\frac{1}{n}\sum_{i=1}^{n}k_{i},$$
con uguaglianza se e solo se i $k_{i}$ sono tutti uguali. Ricordando che $k_{1}+\ldots+k_{n}=120$, otteniamo pertanto che
\begin{eqnarray}
	\frac{9}{2}\sum_{i=1}^{n}k_{i}^{2}+3000n & \geq & 
	\frac{9}{2}\cdot\frac{1}{n}\left(\sum_{i=1}^{n}k_{i}\right)^{2}+3000n
	\nonumber  \\
	 & = & \frac{9}{2}\cdot\frac{1}{n}\cdot 120^{2}+3000n
	\nonumber  \\
	 & = & 600\left(\frac{108}{n}+5n\right),
	\label{stimatempi3}
\end{eqnarray}
con uguaglianza se e solo se $k_{1}=\ldots=k_{n}=120/n$ (cosa che, dovendo nel nostro caso i $k_{i}$ essere tutti interi, non può
accadere se $n$ non è un divisore di 120). \medskip

\emph{Passo 4}. L'espressione (\ref{stimatempi3}), al variare di $n$ 
tra gli interi positivi, assume minimo per $n=5$. Questo si congettura facilmente provando a sostituire i primi valori di $n$, rappresentati nel diagramma seguente (in cui non abbiamo considerato il fattore moltiplicativo 600: si noti che per $n=4$ l'espressione vale 47 e per $n=5$ vale $46.6$).

\begin{center}
	\psset{xunit=5ex,yunit=0.2ex}
	\pspicture(0,-15)(8,120)
	
	\multido{\n=1+1}{7}{\psline[linecolor=blue,linewidth=1.5\pslinewidth]{-}(!\n\space 108 \n\space div 5 \n\space mul add)
	(!1 \n\space add 108 1 \n\space add div 5 1 \n\space add mul add)
	\psdots[dotstyle=triangle](!\n\space 108 \n\space div 5 \n\space mul add)
	(!1 \n\space add 108 1 \n\space add div 5 1 \n\space add mul add)}
	
	\psaxes[dy=20,Dy=20](0,0)(8,120)
	
	\endpspicture
\end{center}


Per dimostrarlo formalmente, consideriamo la disequazione 
\begin{equation}
	\frac{108}{n}+5n\geq\frac{108}{5}+25,
	\label{stimatempi4}
\end{equation}
la quale si può riscrivere nella forma
$$\left(5n-\frac{108}{5}\right)(n-5)\geq 0,$$
da cui segue immediatamente che è verificata per ogni $n\geq 5$ e per ogni $n\leq 108/25=4.32$, quindi per ogni valore intero positivo
di $n$. \medskip

\emph{Passo 5}.  Mettendo insieme la (\ref{stimatempi2}), la
(\ref{stimatempi3}) e la (\ref{stimatempi4}) abbiamo dimostrato che
\begin{eqnarray*}
	\mbox{Tempo totale} & \geq & 
	\left(T-\frac{9}{2}\right)\cdot 120-3000+
	\frac{9}{2}\sum_{i=1}^{n}k_{i}^{2}+3000n  \\
	 & \geq & \left(T-\frac{9}{2}\right)\cdot 120-3000+
	 600\left(\frac{108}{n}+5n\right)	 \\
	 & \geq & \left(T-\frac{9}{2}\right)\cdot 120-3000+
	 600\left(\frac{108}{5}+25\right).
\end{eqnarray*}

Per avere le varie uguaglianze deve essere $n=5$, $k_{1}=\ldots=k_{5}=120/5=24$, e dobbiamo arrivare con la vettura scarica ad ogni pit-stop.

La strategia ottimale prevede pertanto di suddividere la gara in 5 tratti di ugual lunghezza, effettuando quindi 4 pit-stop al termine
dei giri 24, 48, 72, 96.}

\renewcommand{\SoluzioniAlternative}{
}

\renewcommand{\Osservazioni}{\TitoloOsservazione Per dimostare la tesi
del Passo 4, potevamo anche ricorrere al calcolo differenziale. Infatti, da uno studio della funzione
$$f(x)=\frac{108}{x}+5x$$
si vede che essa è decrescente nell'intervallo $0<x\leq\sqrt{108/5}$ e crescente per $x\geq \sqrt{108/5}$.
\begin{center}
	\psset{xunit=5ex,yunit=0.2ex}
	\pspicture(0,-20)(8,125)
	
	\psplot[linewidth=1.5\pslinewidth]{1}{8}{108 x div 5 x mul add}
	
	\psaxes[dy=20,Dy=20](0,0)(8,120)
	
	\psset{linecolor=blue, linestyle=dashed}
	\psline(!21.6 sqrt 0)(!21.6 sqrt 108 21.6 sqrt div 5 21.6 sqrt mul add)
	\psline(!0 108 21.6 sqrt div 5 21.6 sqrt mul add)(!21.6 sqrt 108 21.6 sqrt div 5 21.6 sqrt mul add)
	
	\endpspicture
\end{center}

Poiché $4<\sqrt{108/5}<5$, il minimo di tale funzione sugli interi positivi sarà assunto in 4 o 5: non resta quindi che sostituire questi due valori. Vale la pena di sottolineare che è comunque necessario sostituire entrambi i valori, e \emph{non} basta osservare
che $\sqrt{108/5}$ è più vicino a 5 che a 4. }

\renewcommand{\Commenti}{
}